\section{有效性威胁}

本文的主要有效性威胁体现在以下五个方面。

\begin{itemize}
  \item \textbf{外部机制基线的实现仍属文献启发，而非严格复现。} 本文引入的 ReAct-style single-agent baseline 和 execution-feedback baseline 用于覆盖代表性机制路线，但它们并不声称逐字复现某篇论文的完整系统。因此，外部比较更适合支撑“机制层面对比”，而不是“逐篇文献胜负”。
  \item \textbf{子领域稳定性不等于强跨领域泛化。} \texttt{transfer\_iot\_auto} 与主数据集同属专利数据，只能说明系统在相邻技术子领域上的稳定性，不能直接支撑“跨领域迁移能力”。真正的跨领域主张必须建立在 Bugzilla 等异质数据上的独立验证之上。
  \item \textbf{数据集与问题覆盖仍有限。} 即便增强版路线把主问题扩展到五类，并补充软件工程跨领域小验证，任务集合与跨领域样本规模仍然有限。未来若面向更宽泛的软件工程或科学分析场景推广，仍需在更多数据结构、更多任务类型和更强异质性数据集上继续验证。
  \item \textbf{语义评价仍存在模型偏好与人工样本规模限制。} LLM-as-Judge 能够提供稳定的批量比较，但其评分标准仍可能受模型偏好影响\citep{zheng2023judging,ji2023survey}；人工盲评则更接近人类评审视角，但样本规模、评审背景和时间成本会限制其统计力度。因此，两类评价只能相互补充，不能互相替代。
  \item \textbf{知识约束强度尚未实现自适应调节。} 当前实验已经表明知识图谱收益具有任务依赖性，但本文尚未实现根据问题开放度、证据密度和方法可行性自动调节图谱约束强度的策略。这意味着当前框架虽然能够识别“何时可能过约束”，但还不能自动完成更细粒度的强弱切换。
\end{itemize}

\IfFileExists{generated/table_external_llm.tex}{
外部机制基线结果已经被纳入正文后，第一类威胁由“缺少外部对比”收缩为“外部对比并非严格文献复现”；若当前编译版本尚未显示外部表图，则该缺失本身仍构成投稿前必须显式说明的证据限制。
}{
若当前编译版本尚未显示外部机制基线、子领域稳定性、跨领域验证或人工盲评表图，则这些证据池仍处于回填阶段；这种缺失会直接削弱“外部对比”和“外部有效性”的说服力，因此不能被模糊表述为已经完成。
}

这些威胁不会推翻当前已经观察到的机制趋势，但它们要求本文在写作上始终保持口径克制：内部主张只由已冻结的主证据支持，外部结论只在外部证据真正落地后写入正文，人类评价只在真实盲评分数回收后进入最终结论。
